\section{Encoded frames and changes of ownership}\label{sec:frames}

We next construct the elementary operations that will distribute the ground
vector among the parties of Theorem~\ref{thm:compression}.  These are
contractions between specified register layouts.  A change of layout is part
of the map to be constructed: a raw register cannot be transferred to a
different party merely by changing its name.  Private operations, in contrast,
may act on all the registers of one party, irrespective of their spatial
positions.

\subsection{Sheets and adaptive encodings}

For a physical set $A\subseteq\Lambda$, write
$\mathcal H_A=\bigotimes_{x\in A}\mathbb C^q$.  A \emph{sheet} has one raw
register $\mathbb C^q$ at each $x\in\Lambda$ and a specified owner for each
register.  Different sheets have distinct registers.  Their canonical raw
reference vectors are independent copies of $\Omega$, with sites indexed by
position.  Thus ownership does not permute the positions in the definition
of the reference vector.

Fix a constant $B_0>0$.  In this section all applications of
Proposition~\ref{prop:patch} use accuracy
\begin{equation}\label{eq:frame-epsilon}
 \varepsilon=L^{-60}.
\end{equation}
For a center $c\in\mathbb R^2$ and $0<h\le B_0\log L$, let
\[
 Q^-(c,h)=\{x:\abs{x-c}_\infty\le h\},\qquad
 Q^+(c,h)=\{x:\abs{x-c}_\infty\le 2h\}
\]
be the inner and outer squares.  Their register sets are their intersections
with $\Lambda$.  Choose the projector from Proposition~\ref{prop:patch} in
the form
\begin{equation}\label{eq:hole-projector}
 P_{c,h}
   =\sum_{j=1}^{m}\sum_{\ell=1}^{d_j}
       \bigl(|v_{j\ell}\rangle\langle v_{j\ell}|\bigr)_{D_j}
           \otimes\id_{\Lambda\setminus D_j},
 \qquad
 \norm{(\id-P_{c,h})\Omega}\le\varepsilon.
\end{equation}
Here the $D_j$ are nested square samples, every one containing
$Q^-(c,h)\cap\Lambda$ and contained in $Q^+(c,h)\cap\Lambda$; the vectors
$v_{j\ell}\in\mathcal H_{D_j}$ are normalized; the displayed cylinder
ranges are mutually orthogonal; and $\sum_jd_j$ is bounded by one fixed
power of $L$.  The exponent may depend on $B_0$ and the fixed physical
parameters.  Fix these choices once for each center and scale used in the
construction, and use the same choices on different sheets.

Let $\mathcal J_{c,h}$ be a tag register with orthonormal basis
$\{|j,\ell\rangle\}$.  With a fixed one-site unit vector $|0\rangle$, define
\begin{equation}\label{eq:hole-encoder}
 K_{c,h}
  =\sum_{j,\ell}|j,\ell\rangle_{\mathcal J_{c,h}}\otimes
       \left[
       \bigl(|0\rangle^{\otimes D_j}\langle v_{j\ell}|\bigr)_{D_j}
           \otimes\id_{\Lambda\setminus D_j}
       \right].
\end{equation}
The notation $|0\rangle^{\otimes D_j}$ denotes the product zero vector on
the sites of $D_j$.  Orthogonality of the tag basis gives the exact identity
\begin{equation}\label{eq:encoder-contraction}
 K_{c,h}^{\dagger}K_{c,h}=P_{c,h},\qquad \norm{K_{c,h}}\le1.
\end{equation}
In particular, the encoding retains raw registers at every site, even inside
the hole.  Its tag records which radius and which inside vector were used;
it does not assert a low-rank approximation at one common radius.  If the
outer square has empty physical sample, we use the identity encoding with
a one-dimensional tag.

\begin{definition}[Encoded frame]\label{def:frame}
An encoded frame $\mathcal F$ consists of a raw ownership assignment on one
sheet, a collection of holes $(c,h)$ with pairwise disjoint outer squares,
and a specified party holding each hole's tag.  Its encoding $K_{\mathcal F}$
is the product of the maps \eqref{eq:hole-encoder}, with a fixed ordering of
the newly appended tags, and its reference vector is
\begin{equation}\label{eq:frame-vector}
  \Omega_{\mathcal F}=K_{\mathcal F}\Omega.
\end{equation}
A frame without holes is a raw sheet.  An ownership guide in the plane
specifies such an assignment by sampling its labels at the physical sites.
\end{definition}

Encodings of distinct holes have disjoint raw supports and separate tags,
so they commute after the canonical identification of tag orderings.
If the frame has $r$ holes, their projectors also commute, and therefore
\begin{equation}\label{eq:frame-norm}
 \norm{\Omega_{\mathcal F}}
   =\left\|\prod_{a=1}^{r}P_a\Omega\right\|,
 \qquad
 1-r\varepsilon\le\norm{\Omega_{\mathcal F}}\le1.
\end{equation}
For the lower bound, expand $\prod_aP_a-\id$ by successive factors and use
contractivity and \eqref{eq:hole-projector}.  We always compare vectors in
their canonical site and sheet coordinates while separately specifying the
owners of their input and output registers.  Auxiliary sheets may eventually
be retained as private discard registers.

We shall call a change \emph{bounded} if it uses a fixed number of parties
and gates of the type admitted by Theorem~\ref{thm:compression}.  The number
of normalized pair resources per monomial is fixed; their Hilbert-space
dimensions need not be bounded.  For every fixed $a>0$, an operator-norm
approximation of error $O(L^{-a})$ is allowed, with a polynomial expansion
whose exponent can depend on $a$.  Scaling such an approximation by a
factor $1+O(L^{-a})$ restores contractivity.

\subsection{A tensor estimate for bounded patch networks}

The following elementary tensor estimate permits us to replace a complicated
bounded patch network by pair resources without a dependence on the
dimensions of its open wires.  A normalized vector shared by three or more
parties is not itself one of the permitted pair resources.  We therefore
need products of vectors whose factors each belong to at most two parties.
The estimate below first produces a controlled number of product terms,
keeping all open legs of each original tensor together as one group.
In the patch rewrite that follows, we will verify that each such group
belongs to at most two owners.  These groups are determined by the
individual input bras and output kets, rather than by the parties.

For a matrix $M$, write $\norm{M}_1$ for its
sum of singular values and $\norm{M}_2$ for its Hilbert--Schmidt norm.  The
Hilbert-space norm of a tensor equals the Hilbert--Schmidt norm of every
matrix reshaping of that tensor.

\begin{lemma}[Whole-group tensor bound]\label{lem:group-tensor}
Let a finite graph carry at each vertex a Hilbert-space tensor of norm at
most one.  Contract its internal edges by the usual coordinate pairing,
and suppose there are no edges from a vertex to itself.  Group the open
legs according to their original vertices.  The resulting tensor $Z$ has
norm at most one and has matrix nuclear norm at most one across every
bipartition into whole open-leg groups.

If there are $g$ nonempty open-leg groups, then, for every integer $k\ge1$,
there are subspaces of dimension at most $k$ in the respective group spaces
whose tensor-product projection sends $Z$ to a tensor $Z_k$ with
\begin{equation}\label{eq:group-truncation}
 \norm{Z-Z_k}\le\frac{g}{\sqrt{k}}.
\end{equation}
In orthonormal bases of those subspaces, $Z_k$ has at most $k^g$ product
terms, with coefficients of absolute value at most one.
\end{lemma}

\begin{proof}
Merge two vertices, contracting all edges between them at the same time.
After grouping their common indices, the new tensor is a matrix product
up to a transpose, and
\[
  \norm{AB}_2\le\norm{A}_2\norm{B}_2.
\]
If there are no common edges, use a tensor product instead.  Continue by
merging components and always contracting all edges between the two
components being merged.  No internal self-contraction is left behind.
This proves the norm bound, and proves the same bound for every partially
contracted subgraph with its remaining legs left open.  It applies equally
to graphs with cycles.

For a bipartition of the whole open-leg groups, assign vertices without
open legs to either side.  Contract each side internally.  The two tensors
so obtained have norms at most one.  If all cross edges are grouped into
one index, the desired flattening is $AB^{\mathsf T}$.  Hence
\begin{equation}\label{eq:group-nuclear}
  \norm{AB^{\mathsf T}}_1\le\norm{A}_2\norm{B}_2\le1.
\end{equation}
No dimension enters this estimate.

Flatten $Z$ across one open-leg group and its complement.  Let
$s_1\ge s_2\ge\cdots\ge0$ be its singular values.  By
\eqref{eq:group-nuclear}, $\sum_js_j\le1$, and consequently
\[
 \sum_{j>k}s_j^2
     \le s_{k+1}\sum_{j>k}s_j\le\frac{1}{k+1}.
\]
Projection onto its first $k$ left singular directions thus has error at
most $k^{-1/2}$.  Do this for every original open-leg group.  The projections
act on disjoint tensor factors, so telescoping their product gives
\eqref{eq:group-truncation}.  The resulting tensor has norm at most one;
every coefficient in any product orthonormal basis therefore has absolute
value at most one.  Empty or one-group cases have the same conclusion,
with the natural scalar or one-factor interpretation.
\end{proof}

\subsection{Changing a frame on small patches}

For a frame $\mathcal F$, let $H^-(\mathcal F)$ and $H^+(\mathcal F)$ be the
unions of the physical samples of its inner and outer hole squares.  A tag
is unchanged only if its hole, its encoding, and its owner are unchanged.

\begin{lemma}[Small-patch rewrite]\label{lem:small-rewrite}
Consider old and new one-sheet frames.  Only a bounded number of their holes
are affected; all other holes, encodings, and tag owners agree.  Suppose
there are $m=O(1)$ additional square patches $(c_i,u_i)$, with
$0<u_i\le B_0\log L$, satisfying the following conditions.
\begin{enumerate}
\item Every patch's outer square avoids the physical outer footprints of
the untouched holes.  Affected old and new holes also avoid those
untouched footprints.
\item The union of the patch inner squares contains every site whose raw
owner changes and the entire physical outer footprints of every affected
old or new hole.
\item In each patch's outer square, sites outside $H^-(\mathcal F_{\rm old})$
have at most two old owners.  Separately, sites outside
$H^-(\mathcal F_{\rm new})$ have at most two new owners.
\item All raw owners on the involved patch and affected-hole footprints,
and all affected tag owners, belong to a specified bounded list of parties.
\end{enumerate}
The patches may overlap, and an old affected hole may overlap new affected
holes.  There is a canonical contraction $M$ between the two layouts with
\begin{equation}\label{eq:rewrite-reference-error}
 \norm{M\Omega_{\mathcal F_{\rm old}}-
                   \Omega_{\mathcal F_{\rm new}}}
       \le(m+r_{\rm old})\varepsilon,
\end{equation}
where $r_{\rm old}$ is the number of affected old holes.  For every fixed
$a>0$, $M$ admits a contractive approximation $M_a$ with
$\norm{M_a-M}\le L^{-a}$ and a polynomial sum expansion of the kind required
by Theorem~\ref{thm:compression}, using only the specified parties.  In
particular the rewrite is a bounded change with reference-vector error
$O(L^{-20})$.
\end{lemma}

\begin{proof}
Write $K_{\rm old}$ and $K_{\rm new}$ for the products of just the affected
hole encodings, and $K_0$ for the common untouched encodings.  For the
additional patches, let $P_1,\ldots,P_m$ be the projectors of
Proposition~\ref{prop:patch}, again with error $\varepsilon$.  Fix any order
of the patches, and set
\begin{equation}\label{eq:rewrite-map}
 M=K_{\rm new}P_m\cdots P_1K_{\rm old}^{\dagger}.
\end{equation}
This formula uses canonical raw indices between its factors and interprets
the final raw and tag indices with their new owners.  It is the identity on
all other registers.  Every factor is contractive.

The untouched encodings commute through the displayed factors, by the
first condition.  Moreover,
$K_{\rm old}^{\dagger}K_{\rm old}$ is the product of the affected old-hole
projectors.  For arbitrary contractions $A_1,\ldots,A_s$, the elementary
identity
\[
 A_s\cdots A_1-\id
  =\sum_{j=1}^s A_s\cdots A_{j+1}(A_j-\id)
\]
holds; its first summand is $A_s\cdots A_2(A_1-\id)$ and its last is
$A_s-\id$.  Thus
\[
 \norm{(A_s\cdots A_1-\id)\Omega}
      \le\sum_{j=1}^s\norm{(A_j-\id)\Omega}.
\]
Apply this to all the patch and old-hole projectors in the order in which
they occur.  Multiplication by $K_0$ and $K_{\rm new}$ can only decrease
the norm, giving \eqref{eq:rewrite-reference-error}.  This estimate does not
assume that projectors of overlapping patches commute.

We now construct the required sum expansion of \eqref{eq:rewrite-map}.
Expand each tag and each patch projector using
\eqref{eq:hole-projector}--\eqref{eq:hole-encoder}.  There are polynomially
many complete branch choices, because the number of factors is bounded.
Fix one such branch.  On each selected old-hole square the branch of
$K_{\rm old}^{\dagger}$ consumes the input tag by its basis bra, consumes
the raw inputs by product zero-state bras, and supplies a decoded normalized
ket $v_{j\ell}$.  On a new-hole square the branch of $K_{\rm new}$ contracts
against a normalized bra $\langle v_{j\ell}|$ and supplies product zeros
and a fixed tag.  A selected patch term has a normalized bra and a
normalized ket on its selected square; keep these as \emph{two distinct
tensor vertices}.

Factor out the fixed tag bras and kets, the one-site zero bras and kets, and
every raw wire that continues directly from an input to an output without
encountering any selected action.  The first two sorts of factors are
private norm-one maps.  Every direct wire is private as well: a changed
site belongs to some patch inner square, so it belongs to that patch's
selected square in every radius branch and cannot be such a wire.  The
old and new owner of a remaining direct wire are therefore equal.  In
particular, no factor $\id_{\mathcal H_A}$ with a large Hilbert--Schmidt
norm needs to enter a tensor norm estimate.

Follow the other site wires in their chronological order through the
selected actions.  Each one connects a normalized ket vertex to a later
normalized bra vertex, or gives a raw open input or output.  Actions that
do not contain the site simply let its wire continue.  This leaves a graph
of boundedly many normalized tensors and no self-contractions.  We claim
that each nonempty group of open legs at an original vertex is either
entirely input or entirely output, and belongs to at most two parties.

First, every leg of an old-hole decoded ket encounters a later patch bra.
Indeed its selected square lies in that hole's outer footprint, the whole
of which is covered by patch inner squares.  Thus this decoded ket has no
open output legs.  Dually, every leg of a new-hole encoding bra encounters
an earlier patch ket, so that bra has no open input legs.  These assertions
hold for every choice of the old and new radii, including unequal radii
and overlaps between old and new holes.

The only remaining vertices with open legs are patch bras and patch kets.
A patch bra can have only raw input legs, and a patch ket can have only raw
output legs.  An input site inside an old inner hole is always consumed by
the old decoder's zero-state bra, since every selected hole square contains
its inner square.  It is therefore not an open input of a patch bra.
Consequently all open inputs of a patch bra lie in that patch's outer
square outside the old inner holes, and the third condition limits their
old owners to two.  The corresponding output argument uses the new
encoder's zero-state kets and gives at most two new owners for the open
legs of each patch ket.  Intersections of selected patches only add
internal ket-to-later-bra wires; they do not change these statements about
open legs.  This establishes the claim branch by branch, without any
disjointness assumption on the patches.

Let $Z$ be the remaining open tensor, with groups given by these original
vertices.  Lemma~\ref{lem:group-tensor} supplies an approximation $Z_k$ with
error $g/\sqrt{k}$, where $g=O(1)$, and a sum of at most $k^g$ product terms.
An output group vector in such a term is a normalized state on one party
or a normalized bipartite state on its two owners.  An input group vector,
reshaped with the conjugate coordinates as appropriate, is a normalized
private or bipartite covector.  Their order can be chosen by first
contracting all input effects and then appending all output states, since
the groups have disjoint open legs.  After restoring the factored private
maps and direct identities, each product term is precisely a monomial
allowed by Theorem~\ref{thm:compression}.  It uses only boundedly many pair
states and effects, and its scalar coefficient has absolute value at
most one.  A scalar residual tensor is already a coefficient of absolute
value at most one.

For completeness, the Hilbert-space error on $Z$ bounds the operator-norm
error of this branch.  Consume its factored input bras, apply the residual
tensor reshaped as a map, tensor that map with the direct identities, and
append the factored output kets.  All the outer maps are contractions,
and the operator norm of a matrix is at most its Hilbert--Schmidt norm.
Tensoring with a direct identity preserves its operator norm.  Hence the
branch error is at most $g/\sqrt{k}$, with no factor involving the dimensions
of the bypassing wires.

Let $N\le L^{C_1}$ bound the number of branches, absorbing fixed prefactors
by enlarging $C_1$ for $L\ge2$, and let $g\le g_0$.  Summing their errors
gives an operator-norm error at most $Ng_0/\sqrt{k}$.  Choose a polynomial
$k$ so this quantity is at most $\tfrac12L^{-a}$, and call the resulting
sum $\widetilde M$.  Since $M$ is contractive,
$\norm{\widetilde M}\le1+\tfrac12L^{-a}$.  Then
\[
 M_a=\frac{\widetilde M}{1+\tfrac12L^{-a}}
 \quad\hbox{satisfies}\quad
 \norm{M_a}\le1,\qquad \norm{M_a-M}\le L^{-a}.
\]
The total number of monomials is at most $Nk^{g_0}$, their coefficients are
bounded, and their participating parties are the prescribed ones.  Taking,
for example, $a=30$ proves the final assertion.
\end{proof}

\subsection{Splitting a raw buffer and changing a homogeneous owner}

The next two operations use small mutual information.  We state the
purification consequence with its dimension-independent error explicitly.

\begin{lemma}[Splitting consequence]\label{lem:splitting}
Let $T,U,E$ partition $\Lambda$, and let
$b=I_\Omega(T:E)$.  There are finite-dimensional spaces $\mathcal B_T$ and
$\mathcal B_E$, an isometry
\[
 V:\mathcal H_U\longrightarrow\mathcal B_T\otimes\mathcal B_E,
\]
and normalized vectors
$s\in\mathcal H_T\otimes\mathcal B_T$ and
$s'\in\mathcal H_E\otimes\mathcal B_E$ such that
\begin{equation}\label{eq:frame-splitting}
 \norm{(\id_{TE}\otimes V)\Omega-s\otimes s'}
       \le\sqrt{2(1-e^{-b/2})}\le\sqrt b.
\end{equation}
There is no bound on the dimensions of the two new spaces.  In particular,
if $b\le L^{-60}$, the error is at most $L^{-30}$.
\end{lemma}

\begin{proof}
Set $\rho=\rho_{TE}$ and $\sigma=\rho_T\otimes\rho_E$.  Positivity shows
that $\supp\rho\subseteq\supp\sigma$: the expectation of the projector
onto the kernel of either marginal, tensored with identity, is zero, so
that projector annihilates the support of $\rho$.  In particular
$D(\rho\Vert\sigma)=b$ is well defined on this support.

Choose separate purifications $s$ of $\rho_T$ and $s'$ of $\rho_E$.  Enlarge
their purifying factors, if necessary by appending pure ancillas, so that
$\mathcal B_T\otimes\mathcal B_E$ also has room for an isometric image of
the entire $\mathcal H_U$.  Lemma~\ref{lem:fidelity} and
\eqref{eq:fidelity-information} give a purification overlap
\[
 F(\rho,\sigma)\ge e^{-D(\rho\Vert\sigma)/2}=e^{-b/2}.
\]
Its purification characterization supplies an isometry $V$ on the support
used by $\Omega$, attaining this overlap, with real positive phase, with
$s\otimes s'$.  Extra unused dimensions allow extension to an isometry on
all of $\mathcal H_U$; this extension does not change the overlap.  Both
vectors have norm one, so their squared distance is at most
$2(1-e^{-b/2})\le b$.
\end{proof}

The isometry in Lemma~\ref{lem:splitting} will always act privately: all of
$U$, and both of its enlarged output factors, will be at one party.

\begin{lemma}[Homogeneous birth and death]\label{lem:birth}
In an encoded frame, let $T$ be a set of raw sites owned by $P_\circ$ and
disjoint from every outer hole footprint.  The proposed birth changes only
the owner of $T$, from $P_\circ$ to $Q_\circ$.  Keep every hole, its encoding,
and its tag owner fixed.  Define
\begin{equation}\label{eq:birth-partition}
 E=\{x:\text{the old raw owner of }x\ne P_\circ\}
        \cup H^+(\mathcal F),\qquad
 U=\Lambda\setminus(T\cup E).
\end{equation}
If $I_\Omega(T:E)\le L^{-60}$, the birth is a bounded change involving
only $P_\circ,Q_\circ$, with reference-vector error at most $L^{-30}$.
The reversed death obeys the same bound whenever these ownership and
information conditions hold in its surrounding frame, namely the frame
\emph{after} the death.
\end{lemma}

\begin{proof}
Put $\mathcal V=\id_{TE}\otimes V$, with $V$ from
Lemma~\ref{lem:splitting}, and use its normalized vectors $s,s'$.  All of
$T,U$ initially belongs to $P_\circ$.  At that party apply $V$ on $U$,
contract $T\mathcal B_T$ with $\langle s|$, and supply the normalized pair
state $s$ on $T\mathcal B_T$ with $T$ now held by $Q_\circ$ and
$\mathcal B_T$ still held by $P_\circ$.  Finally apply $V^{\dagger}$ to
$\mathcal B_T\mathcal B_E$ at $P_\circ$.  The associated canonical map is
\begin{equation}\label{eq:birth-map}
 B=\mathcal V^{\dagger}
       \left(|s\rangle\langle s|_{T\mathcal B_T}
                            \otimes\id_{E\mathcal B_E}\right)
          \mathcal V.
\end{equation}
It is a Hermitian contraction.  Writing $\Pi$ for the middle projector,
\[
 \norm{(B-\id)\Omega}
   \le\norm{(\Pi-\id)\mathcal V\Omega}
   =\norm{(\Pi-\id)(\mathcal V\Omega-s\otimes s')}
   \le L^{-30}.
\]
Every hole encoder is supported in $E$, on whose raw registers this map
acts as the identity.  It also acts as the identity on tags.  Commuting the
encoders through the map therefore gives the same error bound on the
encoded reference vectors, by their contractivity.

For the reversed operation, first apply $V$ privately on $U$ at
$P_\circ$, use the normalized bipartite effect $\langle s|$ on
$T\mathcal B_T$ held by $Q_\circ,P_\circ$, and reinsert $s$ wholly at
$P_\circ$.  Applying $V^{\dagger}$ gives the same canonical operator
$B=B^{\dagger}$, with the reverse input and output ownership.  This has
one pair effect rather than one pair source.  Both procedures have the
required bounded expansion.  If the two owner labels agree, all these
maps are private; the identity also realizes the ownership change.
\end{proof}

In the geometric construction the information hypothesis is supplied by
Proposition~\ref{prop:angular}, with its exponent set to $60$.  In
particular, for finitely many marked points and a fixed clearance constant,
the logarithmic floor $t=D\log L$ can be chosen so that any eligible
tapering pair $T,E$ satisfies the hypothesis.  The set $E$ in
\eqref{eq:birth-partition} includes outer hole footprints even when all
their raw owners happen to be $P_\circ$; this keeps the splitting isometry
off every encoding support.

\subsection{Exchanging assignments between two sheets}

\begin{lemma}[Two-sheet exchange]\label{lem:exchange}
Consider two encoded frames, and a plane region $Y$.  Suppose every hole's
entire outer square lies on one side of $\partial Y$.  Inside $Y$, exchange
the two raw ownership assignments and the two lists of encoded holes,
including their tag owners; outside $Y$, keep both frames unchanged.
Let $P_\circ$ be a party, and on the physical square put
\begin{align}
 Z&=\{x:\text{the two old raw owners of }x
                         \text{ are not both }P_\circ\}
                  \cup H^+(\mathcal F_1)\cup H^+(\mathcal F_2),
                 \label{eq:exchange-Z}\\
 T&=Z\cap Y,\qquad E=Z\setminus Y,\qquad
 U=\Lambda\setminus Z.\label{eq:exchange-partition}
\end{align}
If $I_\Omega(T:E)\le L^{-60}$, the exchange can be implemented by private
contractions and register renaming, with reference-vector error at most
$4L^{-30}$.  It uses $P_\circ$ and the owners of the registers and tags
being renamed.  Thus, when these parties form a bounded list, it is a
bounded change in the sense of Theorem~\ref{thm:compression}.
\end{lemma}

\begin{proof}
Let $K_{\rm in}$ and $K_{\rm out}$ denote the products of the two frames'
encoders before and after the exchange.  For $A\subseteq\Lambda$, write
$F_A$ for the unitary swapping the two raw sheets at the positions in $A$.
First rename the two sheets inside $Y$, together with each tag of a hole
there, leaving every physical register at its present party.  Call this
map $\mathcal R$.  This is a private identification of tensor factors: the
output owner of sheet 1 at $x\in Y$ is precisely the old owner of sheet 2
there, and conversely.  The same assertion holds for a tag moved with its
whole hole.  In canonical raw coordinates the exact intertwining identity
is
\begin{equation}\label{eq:exchange-intertwining}
  \mathcal R K_{\rm in}=K_{\rm out}F_{Y\cap\Lambda}.
\end{equation}
It follows by checking each site outside holes and each whole hole
encoding separately.  For a hole within $Y$, all selected raw squares and
its tag change sheet together; for a hole outside $Y$, nothing changes.
The assumption on the outer square ensures this for every selected radius.

All registers of $U$ on both sheets are raw, outside every hole, and held
by $P_\circ$ both before and after this renaming.  Apply
Lemma~\ref{lem:splitting} to the one-copy partition $T,U,E$, and again write
$\mathcal V=\id_{TE}\otimes V$.  Let
\begin{equation}\label{eq:exchange-buffer-correction}
 D_U=(V^{\otimes2})^{\dagger}F_{\mathcal B_T}V^{\otimes2}
\end{equation}
on the two copies of $U$, where $F_{\mathcal B_T}$ swaps the two enlarged
$\mathcal B_T$ factors.  The whole map is a private contraction at
$P_\circ$: apply the two copies of $V$, swap the indicated factors, and
apply their adjoints.

Set $A=Y\cap U$.  After $\mathcal R$, first apply the private swap $F_A$,
and then $D_U$.  These two corrections need not commute, so we retain this
order.  Both corrections commute through $K_{\rm out}$, since they act
only on raw sites of $U$, away from every hole, and act trivially on tags.
As $Y\cap\Lambda=T\sqcup A$, their net action on raw two-copy coordinates
is
\begin{align}
 D_UF_AF_{Y\cap\Lambda}
   &=D_UF_T\notag\\
   &=(\mathcal V^{\otimes2})^{\dagger}
          F_{T\mathcal B_T}\mathcal V^{\otimes2}.
              \label{eq:exchange-net-map}
\end{align}
The last equality uses that $F_T$ acts outside $U$ and hence commutes
through its isometries.  Consequently the implemented encoded map
$C=D_UF_A\mathcal R$ satisfies
\begin{equation}\label{eq:exchange-exact-map}
 CK_{\rm in}
   =K_{\rm out}(\mathcal V^{\otimes2})^{\dagger}
                   F_{T\mathcal B_T}\mathcal V^{\otimes2}.
\end{equation}

Put $\zeta=s\otimes s'$ and
$\delta=\norm{\mathcal V\Omega-\zeta}\le L^{-30}$.  Both vectors are
normalized, so
\[
 \norm{(\mathcal V\Omega)^{\otimes2}-\zeta^{\otimes2}}
      \le2\delta.
\]
The vector $\zeta^{\otimes2}$ is fixed by $F_{T\mathcal B_T}$, which swaps
the two identical $s$ factors.  Since
$(\mathcal V^{\otimes2})^{\dagger}\mathcal V^{\otimes2}=\id$, we obtain
\begin{align*}
 &\left\|
   (\mathcal V^{\otimes2})^{\dagger}F_{T\mathcal B_T}
        \mathcal V^{\otimes2}\Omega^{\otimes2}-\Omega^{\otimes2}
   \right\|\\
 &\hspace{2em}\le
   \norm{(F_{T\mathcal B_T}-\id)
                    ((\mathcal V\Omega)^{\otimes2}-\zeta^{\otimes2})}
       \le4\delta.
\end{align*}
Multiplication by the contraction $K_{\rm out}$ and
\eqref{eq:exchange-exact-map} give the claimed error between the product
of the old encoded reference vectors and the product of the new ones.
All actual operations are private at the indicated parties; the apparent
partial raw swap was induced by the change of sheet names and corrected
using the common owner's buffer.  No raw register is transferred by this
renaming.
\end{proof}

Proposition~\ref{prop:angular} will again verify the one-copy information
bound in Lemma~\ref{lem:exchange}.  Including holes from \emph{both} sheets
in \eqref{eq:exchange-Z} ensures that the private correction commutes with
every encoding.  Exchanging whole outer squares ensures that all tag and
radius branches use the same register identification.  The second sheet
may be discarded after the exchange and any subsequent point treatments;
its encoded reference vector is retained in the error comparison until
that discard is taken.
